%=======================================================================
%  MusIML @ NeurIPS 2026 -- Track 5, LONG PAPER (8 pages).
%
%  Venue : NeurIPS.cc/2026/Workshop/MusIML  (Sydney)
%  Due   : 2026-09-09 20:00 UTC (07:00 AEDT 10 Sep), 30-min grace.
%  Rules : 8 pages excluding references; figures and tables COUNT.
%          Anonymised, double-blind.  Non-archival by default.
%
%  REQUIRES neurips_2026.sty in the same directory.  Download from the
%  NeurIPS 2026 Overleaf template:
%    https://www.overleaf.com/latex/templates/
%      formatting-instructions-for-neurips-2026/bjdwqfdkyftc
%
%  This is a condensed version of iciser2026_short.tex.  Every number is
%  transcribed from that file, which in turn comes from
%  notebooks/GRADE_figures.ipynb and results/loto_runs.csv (seed 42).
%=======================================================================
\documentclass{article}

% Anonymised for double-blind review.  Use \usepackage[final]{neurips_2026}
% for a camera-ready with author names.
\usepackage{neurips_2026}

\usepackage[utf8]{inputenc}
\usepackage[T1]{fontenc}
\usepackage{hyperref}
\usepackage{url}
\usepackage{booktabs}
\usepackage{amsfonts}
\usepackage{nicefrac}
\usepackage{microtype}
\usepackage{graphicx}
\usepackage{amsmath}

% Seven floats in eight pages: let LaTeX use bottoms and more of
% each page before it defers a float to a later page.
\renewcommand{\topfraction}{0.95}
\renewcommand{\bottomfraction}{0.85}
\renewcommand{\textfraction}{0.05}
\renewcommand{\floatpagefraction}{0.6}
\setcounter{topnumber}{3}
\setcounter{bottomnumber}{2}
\setcounter{totalnumber}{4}

\title{GRADE: Do AI-Text Detectors Survive the Next Assignment,\\
and Who Pays When They Fail?}

\author{%
  Anonymous Author(s)\\
  Affiliation\\
  Address\\
  \texttt{email}
}

\begin{document}

\maketitle

\begin{abstract}
Generative AI has unsettled written assessment, and universities now screen
student essays with AI-text detectors reporting accuracies above 99\%,
usually cross-validated on the very corpus they are scored against. We
introduce GRADE (Generalization Robustness Assessment for Detectors in
Education), a leave-one-assignment-out evaluation over 44{,}861 student
essays across 15 assignments that attributes every false accusation to the
student who wrote it. Against a word-frequency baseline fitted on every
fold, accuracy on an unseen assignment stays high (93--100\%), yet the
baseline matches or beats \texttt{deberta-v3-small} on 9 of 15 folds and
\texttt{roberta-base} on 8: high accuracy is not evidence of detecting
authorship. Both exceed 99.6\% ROC-AUC everywhere, locating the failure in
the decision threshold, not the representation. False accusations also fall
5.6 times as often on the most-highly-graded students as on the least,
traced to formal style rather than grade. We release the merged,
deduplicated corpus with assignment labels and named generators.
\end{abstract}

\section{Introduction}

A student can now produce a competent essay without doing the thinking the
essay was set to evidence. Institutions are under pressure to respond, and
one response is technical: adopt a detector, run submissions through it,
and treat the output as evidence in academic-integrity processes.

That response is justified by numbers. Recent work reports Macro F1 scores
of 99.78\% and 98.10\% for fine-tuned detectors, rising to 100.00\% on one
corpus \cite{qorich2025}. Commercial detectors report comparably high
figures. Read at face value, the problem looks solved.

It is not solved. Those figures are measured on the corpus the detector was
built from. The essays it is scored against answer the same assignments,
come from the same cohort, and were generated by the same models as the
essays it learned from. Choosing a stronger architecture does not close
this gap, because the number being compared across architectures is itself
measured the same way. A programme director needs two other things no such
number speaks to: what happens with next semester's assignment, which by
definition no detector has been scored against yet, and what happens to the
students the tool gets wrong.

We ask two questions. \textbf{RQ1.} If a detector is trained on a set of
assignments and then meets an assignment absent from training, how much of
its reported performance survives? \textbf{RQ2.} When it is wrong about
authentic student work, which students absorb the error?

RQ1 isolates the one shift a programme cannot avoid: every term brings a
new assignment. RQ2 asks what an aggregate accuracy conceals, because a
false positive is not a statistical event. It is an accusation against a
named student. Our contributions are:

\begin{itemize}\itemsep2pt
\item A leave-one-assignment-out protocol that measures transfer to a
genuinely unseen assignment, rather than to a held-out sample of the same
assignments a detector trained on.
\item A word-frequency control fitted on every fold, which is sufficient on
a majority of folds to match a fine-tuned transformer, so the accuracy a
detector reports is not by itself evidence of detecting authorship. Two
encoders reproduce this, and their ROC-AUC places the loss on the
threshold, not the representation.
\item A join of detector verdicts to teacher-assigned grades and writer
characteristics for the same essays, showing the false-positive rate rises
with grade rather than falling on it, with a formality mechanism tested and
supported.
\end{itemize}

\section{Related Work}

Detecting generated text is usually posed as supervised binary
classification: a pretrained encoder is fine-tuned on a corpus pairing
human-written documents with generated ones, and the resulting classifier
is scored on held-out items from that same corpus. Qorich and El Ouazzani
are representative of the strongest current results, fine-tuning GPT-2,
RoBERTa and BART across five corpora \cite{qorich2025} and measuring
generalization by five-fold cross-validation on each target corpus. Their
numbers are not in dispute, but cross-validation trains on four fifths of
the corpus it evaluates, so a detector that had memorised the vocabulary of
those particular prompts, and learned nothing transferable, would still
score well. The same design constrains what can be said about who a
detector gets wrong: a group difference under 2\% is reported for essays
that carry no demographic field \cite{qorich2025}. Neither limitation is
peculiar to that study. Both follow from an evaluation protocol the field
has largely settled on. The same pattern holds in work that takes
interpretability seriously: a DeBERTa-v3 detector reporting 98.21\%
validation accuracy adds token-level explanations to make its decisions
auditable \cite{sentinel2026}, and a retrieval-augmented detector pairs an
explanation layer with defences against obsolescence and paraphrasing
attacks \cite{rag2026}. Both are scored in distribution, so neither speaks
to what survives an assignment the detector has never seen.

Work situated in education reaches conclusions difficult to reconcile with
accuracies above 99\%. Weber-Wulff et al. evaluate fourteen tools,
including commercial ones marketed to universities, and find none reliable
enough for academic-integrity decisions, degrading further under
obfuscation \cite{weberwulff2023}. Sadasivan et al. give theoretical
grounds suggesting reliable detection may not be achievable as generated
and human text converge \cite{sadasivan2023}. Liang et al. matter most to
RQ2: they show detectors flag non-native English writing as generated far
more often than native writing \cite{liang2023}. That finding is bounded by
its design, because the non-native and native writing come from different
sources, so writer group and corpus vary together. Whether the same gap
appears between writers inside one cohort, answering the same assignment
and marked by the same raters, is open.

Robustness under distribution shift is an established concern, and the
benchmarks built to measure it are explicit about which shift they vary.
RAID spans eleven generators, eight domains and eleven adversarial attacks,
and finds detectors degrade sharply outside their tuning conditions
\cite{raid2024}. MAGE reports substantial degradation on unseen generators
\cite{mage2024}, and paraphrasing is a documented evasion strategy
\cite{krishna2023}. These vary the generator, the genre and the provenance.
None of those is what changes inside a running course: the cohort, the
marking scheme and the tools stay largely fixed from one term to the next,
while the assignment is replaced. We are not aware of a study that isolates
that shift, nor of one reporting which students absorb a detector's false
accusations. GRADE addresses both.

\section{Data}

Answering RQ1 and RQ2 requires a corpus of student essays large enough for
15 held-out folds, labelled with the assignment each essay answers, and
joinable to real writer information. Both source collections draw their
student essays from PERSUADE: 25,996 argumentative essays written by United
States students in grades 6 to 12 across 15 assignments \cite{crossley2024}.
The first is the corpus released for the LLM-Detect competition
\cite{king2023}; the second is a community-assembled corpus for the same
task, pairing PERSUADE essays with generated essays from 16 named models
including Llama-2, Mistral, Claude, PaLM and Falcon \cite{daigt2024}.

Only the community corpus records which assignment each essay answers, and
65\% of the competition corpus is already contained within it, so every
labelled essay comes from the community corpus: \textbf{44,861 essays across
15 assignments}, 27,367 student to 17,494 generated. We match the two on
normalised text rather than stacking them, because stacking would place the
same essay in the corpus twice and let a later split put one copy in
training and the other in test. Two essays carried conflicting labels and
were discarded; 9,481 competition essays that no assignment label covers
are excluded rather than assigned by inference.

PERSUADE additionally records, for each essay, the holistic grade awarded
by human raters on a 1 to 6 scale, the school year, whether English is an
additional language, and indicators of economic disadvantage and disability
\cite{crossley2024}. Matched on exact text, it covers 25,995 of the 27,367
student essays, or 95\%. PERSUADE contains no generated text and is
therefore never trained on: it is a reference layer applied to the
detectors' verdicts, and it is what makes RQ2 answerable.

\section{Method}


The protocol has three parts: a detector fine-tuned the same way on every
fold, a leave-one-assignment-out split that rotates which assignment is
withheld, and an attribution step that joins verdicts back to the writer of
each essay. Figure~\ref{fig:arch} gives an overview before each part is
described in turn.

\paragraph{Detectors.}
We fine-tune two parameter-matched, architecturally distinct encoders
end-to-end for binary classification. \texttt{deberta-v3-small} ($\approx$140M
parameters) uses disentangled attention, decomposing the score between
positions $i$ and $j$ into content-to-content, content-to-position and
position-to-content terms. \texttt{roberta-base} ($\approx$125M) is comparably
sized but computes a single undecomposed attention over content and
position jointly. Both truncate to 256 tokens and pass the contextual
\texttt{[CLS]} representation to a feed-forward head. We run the full
rotation once per encoder, not to rank the two but to establish whether
what the rotation exposes is a property of detectors or of one model.
Training is identical across both encoders and all 15 folds: AdamW,
learning rate $2\times10^{-5}$, batch size 8, 5 epochs, cross-entropy, full
fine-tuning, best checkpoint on the stopping portion, seed 42, giving
15 folds $\times$ 2 encoders $=$ 30 runs. Nothing is tuned per assignment,
so a difference between folds reflects the assignment rather than the
effort spent on it. Table~\ref{tab:recipe} gives the full
configuration.

\begin{table}[t]
\centering
\small
\caption{Training configuration, fixed before the first fold and identical
across both encoders, all 15 folds and the in-fold controls.}
\label{tab:recipe}
\begin{tabular}{@{}ll@{}}
\toprule
\textbf{Setting} & \textbf{Value} \\
\midrule
Conditions        & \texttt{deberta-v3-small}, \texttt{roberta-base} \\
Initialisation    & pretrained checkpoint, fresh head \\
Max token length  & 256 \\
Optimiser         & AdamW \\
Learning rate     & $2\times10^{-5}$ \\
Batch size        & 8 \\
Epochs            & 5 \\
Loss              & cross-entropy \\
Fine-tuning       & full, encoder and classifier \\
Checkpoint        & best on the stopping portion \\
Random seed       & 42 \\
Total runs        & 15 folds $\times$ 2 encoders $=$ 30 \\
\bottomrule
\end{tabular}
\end{table}

\begin{figure}[!ht]
\centering
\includegraphics[width=\textwidth]{Assets/grade_arch.pdf}
\caption{The GRADE protocol. Two collections of the same student writing
are merged and deduplicated into 44{,}861 essays across 15 assignments. One
assignment is held out at a time: both encoders are fine-tuned under one
recipe on the remaining 14 and marked on the one withheld. Three split
properties are enforced by audits that halt the study rather than warn, and
two controls that cannot recognise authorship are refitted on every fold.
Teacher-assigned grades and writer characteristics are then joined to the
verdicts, so each false accusation is attributed to the student who wrote
the essay.}
\label{fig:arch}
\end{figure}

\paragraph{The leave-one-assignment-out protocol.}
The intuition is familiar from masked language modelling, where tokens are
hidden and recovered from surrounding context. GRADE masks at a different
granularity and for the opposite purpose: the unit is an entire assignment
rather than scattered tokens, and recovering any of it from training would
be contamination rather than success. Formally the protocol is grouped
cross-validation with the assignment as the group, in the manner of
leave-one-subject-out designs in clinical machine learning. Let
$\mathcal{A}$ be the set of 15 assignments; for each held-out
$a \in \mathcal{A}$ we set $\mathrm{Train}(a) = \mathcal{A}\setminus\{a\}$
and $\mathrm{Test}(a) = \{a\}$. The 14 assignments in $\mathrm{Train}(a)$
are divided 85\% to fit the detector and 15\% to decide when to stop
fitting it. The withheld assignment contributes no essay to either portion
and is opened once, at the end. Cohort, marking scheme and generators are
held fixed across folds; only the assignment differs.

The excluded essays matter here because any of them could belong to the
assignment being withheld. A topic classifier accurate to 97.4\% on
known assignments places 62\% of them with no confidence, indicating
sources beyond the 15 and confirming they cannot be assigned safely.

\paragraph{What the split guarantees.}
This corpus is merged from two collections that overlap, so the same essay
text can arrive from more than one source row. We therefore enforce three
properties on every fold rather than assuming them. No essay from the
held-out assignment appears in training or in the stopping portion. No text
appears in more than one of the three portions. And every portion carries
equal numbers of student-written and generated essays, which fixes a
guessing detector at 50\% and makes folds comparable. Each property is
checked by an audit that raises and halts the study rather than issuing a
warning, so a violation cannot pass unnoticed into a reported number. All
15 folds pass all three checks.

\paragraph{Controls.}
Accuracy alone cannot say whether a detector has learned anything about
authorship, because a method that plainly has not may reach the same
number. Alongside every detector we fit two deliberately unintelligent
comparisons, each refitted on that fold's training portion. The first sees
a single number per essay: its length. The second sees which words appear
and how often, with no sense of order or meaning. Neither has any notion of
who wrote a text, so whatever they score is reachable without recognising
authorship. A detector must beat them before its accuracy can be credited
to anything about writing. Length is a low bar here by construction: the
median student essay is 2,131 characters against 1,974 generated. Word
frequency asks the harder question, and is refitted per fold because that
length balance does not hold on every assignment.

\paragraph{Attributing the errors.}
For RQ2 we take every authentic essay in a held-out assignment, join the
detector's verdict to what PERSUADE records about its writer, and report
the false-positive rate per group. Only authentic work enters this
analysis, because only authentic work can be falsely accused. The join is
on the essay identifier rather than on row order.

\section{Results}

Generated text is the positive class. We report accuracy, recall and
ROC-AUC, and both error rates where they carry the argument: the
false-positive rate is the share of genuine student essays wrongly flagged,
the false-negative rate the share of generated essays passed as authentic.
Training is stochastic, so we rest the argument on aggregates over the 30
runs and treat the ordering of adjacent folds as provisional.

\subsection{RQ1: transfer to an unseen assignment}

The detector transfers. Accuracy ranges from 93.47\% (Exploring Venus) to
99.61\% (Car-free cities), and it clears the length-only control on every
fold, by 13 to 50 points and 34.0 on average. Read alone, this is the
number a vendor would quote: an unseen assignment costs a programme almost
nothing.

Read against the word-frequency control, it is a different result. The
detector loses to that control on 9 of the 15 folds, by a mean of 0.76
points and by as much as 5.25 on Exploring Venus, its own worst fold. A
classifier with no notion of authorship, built from which words appear and
how often, matches or beats a fine-tuned transformer more often than not
under the exact shift this paper studies. Performance survives the shift
almost completely, but what survives is not shown to be detection of
authorship rather than of vocabulary. Table~\ref{tab:loto} reports every fold, and
Figure~\ref{fig:margins} plots the margin against that control for both
encoders.

\begin{table}[tb]
\centering
\small
\caption{Leave-one-assignment-out, all 15 folds, seed 42, identical
recipe for both encoders. Each row trains on the other 14 assignments and
marks only the one named; generated text is the positive class and all
figures are percentages. $\Delta$ is accuracy minus the word-frequency
control fitted on that same fold. Both encoders rank almost perfectly on
every assignment, and both lose accuracy where the decision threshold does
not transfer.}
\label{tab:loto}
\begin{tabular}{@{}lrrrrrrr@{}}
\toprule
& & \multicolumn{2}{c}{\textbf{DeBERTa-v3}} & \multicolumn{2}{c}{\textbf{RoBERTa}}
& \multicolumn{2}{c}{\textbf{$\Delta$ vs control}} \\
\cmidrule(lr){3-4}\cmidrule(lr){5-6}\cmidrule(lr){7-8}
\textbf{Held-out assignment} & \textbf{Test $n$} & \textbf{Acc} & \textbf{AUC} & \textbf{Acc} & \textbf{AUC} & \textbf{DeB} & \textbf{RoB} \\
\midrule
Distance learning            & 4{,}314 & 98.79 & 99.85 & 96.64 & 99.80 & $-$0.65 & $-$2.80 \\
Car-free cities              & 4{,}102 & 99.61 & 99.99 & 99.59 & 99.99 & +0.29 & +0.27 \\
Electoral college            & 3{,}440 & 99.48 & 100.00 & 90.06 & 99.85 & +0.53 & $-$8.89 \\
Seeking multiple opinions    & 3{,}104 & 96.01 & 99.68 & 98.16 & 99.97 & $-$2.19 & $-$0.04 \\
Mandatory extracurriculars   & 2{,}812 & 98.36 & 99.82 & 98.12 & 99.68 & $-$0.64 & $-$0.88 \\
Summer projects              & 1{,}902 & 97.79 & 99.81 & 97.21 & 99.77 & $-$1.37 & $-$1.95 \\
Facial action coding         & 1{,}834 & 99.24 & 99.97 & 99.84 & 99.99 & +0.77 & +1.37 \\
Community service            & 1{,}100 & 98.64 & 99.97 & 99.82 & 100.00 & +0.09 & +1.27 \\
A Cowboy Who Rode the Waves  & 1{,}046 & 96.46 & 99.96 & 98.47 & 99.97 & $-$3.16 & $-$1.15 \\
Grades for extracurriculars  &     980 & 99.18 & 99.94 & 99.08 & 99.99 & $-$0.41 & $-$0.51 \\
Cell phones at school        &     926 & 99.46 & 99.97 & 99.35 & 99.99 & $-$0.11 & $-$0.22 \\
Phones and driving           &     830 & 98.67 & 99.97 & 97.71 & 99.94 & +1.07 & +0.11 \\
Driverless cars              &     728 & 99.59 & 99.99 & 99.45 & 100.00 & +0.69 & +0.55 \\
Exploring Venus              &     628 & 93.47 & 99.98 & 99.84 & 100.00 & $-$5.26 & +1.11 \\
The Face on Mars             &     620 & 96.94 & 99.70 & 99.35 & 99.99 & $-$1.12 & +1.29 \\
\midrule
\textbf{Mean} & & 98.11 & 99.91 & 98.18 & 99.93 & $-$0.76 & $-$0.70 \\
\textbf{Folds above control} & & & & & & 6/15 & 7/15 \\
\midrule
\multicolumn{8}{@{}l}{\emph{controls, mean over folds: word-frequency 98.88, length-only 64.16}} \\
\bottomrule
\end{tabular}
\end{table}

\begin{figure}[tb]
\centering
\includegraphics[width=\textwidth]{Assets/rq1_margin_vs_control.pdf}
\caption{Accuracy minus the word-frequency control on each held-out
assignment, both encoders. Bars to the left of zero are assignments where a
method that cannot recognise authorship scored at least as well as the
fine-tuned detector.}
\label{fig:margins}
\end{figure}

\subsection{The result is not an artefact of one encoder}

Repeating all 15 folds with \texttt{roberta-base} leaves the aggregate
picture unchanged. The two encoders are indistinguishable on every summary
measure in Table~\ref{tab:loto}, and RoBERTa loses to the
word-frequency control on 8 folds against DeBERTa's 9. Changing the encoder
does not recover the margin, so the finding belongs to the evaluation
rather than the model. Where they differ is more informative: each one's
worst assignment is close to the other's best, and their margins correlate
at $-0.15$ across the 15 folds, which is to say not at all. One detector's
weak assignment is no warning for another's.

Neither collapse is representational. ROC-AUC never falls below 99.68\% for
either encoder, so both rank essays correctly on every assignment; on its
worst fold RoBERTa still recovers only 80.17\% of generated essays at the
default cut point. What fails is the decision threshold, fitted where the
assignment was seen and not travelling with the ranking.

Re-scoring each fold at the cut point that maximises accuracy on it bounds
what better calibration could buy (Figure~\ref{fig:threshold}). The two
worst folds recover to 98.84\% and 97.61\%, and they need opposite
corrections, 0.025 against 0.985, so no single fixed threshold would have
served both. No detector can perform the correction anyway, because
choosing that cut point requires the labels for the assignment being
marked. The bound is the important part: given the threshold it could not
have known, each encoder still clears the word-frequency control on only 9
of the 15 folds, by 0.21 and 0.46 points on average. Calibration explains
where the accuracy went; it does not supply the evidence of detecting
authorship that the control was there to test for.


\begin{figure}[tb]
\centering
\includegraphics[width=0.62\textwidth]{Assets/threshold_headroom.pdf}
\caption{Accuracy at the default cut point, with the shaded extension
showing what a threshold chosen on the held-out assignment itself would
recover. The correction is unavailable in use: it needs the labels the
detector is being asked to supply.}
\label{fig:threshold}
\end{figure}

\subsection{RQ2: who absorbs the error}

The false-positive rate rises with grade at every step, from 0.92\% at
grade 1 to 5.18\% at grade 6, a 5.6-fold increase with no reversal across
the six bands (Table~\ref{tab:fairness}). English-language status shows the
same shape in the opposite direction to prior work: writers for whom
English is an additional language are flagged at 0.63\%, against 2.87\% for
other writers, a 4.6-fold difference. Prior work finds non-native writers
over-flagged; here, inside one corpus with grade and rater held fixed, they
are under-flagged. The two results answer RQ2 the same way: the detector's
few mistakes are not spread evenly across students, and they concentrate on
writers whose English is judged the strongest by an independent human
rater, not the weakest. The comparison with Liang et al. \cite{liang2023}
is cleaner in one respect: both writer groups sit inside the same corpus
and are marked by the same raters, so group and corpus do not vary together
as in their design.

\begin{table}[tb]
\centering
\small
\caption{The false-positive rate on authentic student work, pooled over the
15 held-out folds. Every figure is an accusation against a student who did
their own work.}
\label{tab:fairness}
\begin{tabular}{@{}lrr@{}}
\toprule
\textbf{Writer group} & \textbf{Essays} & \textbf{False-positive rate} \\
\midrule
\multicolumn{3}{@{}l}{\emph{by holistic grade}} \\
\quad 1 (weakest)   & 434   & 0.92 \\
\quad 2             & 2{,}314 & 1.30 \\
\quad 3             & 3{,}805 & 1.47 \\
\quad 4             & 3{,}846 & 3.35 \\
\quad 5             & 2{,}166 & 4.29 \\
\quad 6 (strongest) & 637   & 5.18 \\
\midrule
\multicolumn{3}{@{}l}{\emph{by English-language status}} \\
\quad English an additional language     & 1{,}271  & 0.63 \\
\quad English not an additional language & 11{,}483 & 2.87 \\
\bottomrule
\end{tabular}
\end{table}

\subsection{Why: testing formality}

A rate that rises with grade is a correlation, not yet a mechanism, so we
test one directly. Generated essays in this corpus carry formal discourse
markers at roughly four times the rate of student essays, 4.98 per thousand
words against 1.25, with a median of 3.7 against zero. A detector could
therefore learn formal polish as evidence of generation, and would then
flag the students who write most formally. Testing this means asking
whether marker density separates flagged from passed essays \emph{within} a
grade band, which is what would show polish predicting a flag beyond what
the grade already explains. We use every authentic essay, with no sampling,
and the marker list was fixed from the discourse literature before any
verdict was examined.

Within every grade band with enough flagged essays to compare, flagged
essays carry a higher density of formal discourse markers than passed
essays from the same band: a difference of 1.30 markers per thousand words
at grade 2, rising to 1.77 at grade 5 and holding positive at every band in
between (Figure~\ref{fig:markers}). Grade is therefore not the whole
explanation; formality predicts a flag beyond what the grade accounts for.
A token-attribution check with LIME \cite{ribeiro2016} over 60 flagged
essays surfaced mostly ordinary content words and is uninformative at that
sample size.

The reading we draw is that the detector has partly learned formality as a
proxy for generation: a student who writes formally resembles the generated
signature whether or not the essay is genuine. This also explains the
English-language result. Writers for whom English is an additional language
write less densely in these markers on average, and are flagged less often
as a consequence, not because they write worse.

\begin{figure}[tb]
\centering
\includegraphics[width=0.78\textwidth]{Assets/rq2_marker_density.pdf}
\caption{Formal discourse-marker density in authentic essays the detector
flagged against those it passed, compared within each grade band. The
flagged essays are denser in every band, so formality predicts a flag
beyond what grade accounts for.}
\label{fig:markers}
\end{figure}

\section{Conclusion and Future Work}

We introduced GRADE, a protocol for measuring what an AI-text detector
retains when it meets an assignment absent from training, and for
attributing its false accusations to the students who wrote the work.
Accuracy on an unseen assignment stays high, 93.47\% to 99.61\% across the
15 folds, but a classifier built only from word frequencies matches or
beats the fine-tuned detector on 9 of those 15, so high accuracy does not
by itself show detection of authorship. A second encoder reproduces that
result and locates the failure: ROC-AUC never falls below 99.68\% for
either model, so both rank essays correctly on every assignment and what
does not survive the shift is the decision threshold --- though repairing
that threshold with labels no detector would have still leaves both
encoders clear of the control on only 9 of 15 folds. And the false
positives are not spread evenly: students with the strongest holistic
grades are accused 5.6 times as often as students with the weakest, which a
test within grade bands ties to formal writing style rather than to grade
as such.

On this evidence we propose that a university evaluating a detector require
five things: performance on student work the supplier did not build the
tool from; the score of a word-frequency comparison, not only a
length-only one, on the same evaluation set; the two error rates separately
rather than a single accuracy; performance on an assignment absent from
training; and the false-positive rate broken down by writer group.

These conclusions are bounded. The evidence is one cohort and one corpus
family, scored by two encoder-only detectors at a single seed, so the
ordering of adjacent folds is provisional and the pattern is untested on
university coursework or on much larger, zero-shot and commercial
detectors. The 9,481 essays carrying no assignment label are excluded rather than
assigned by inference, so this study's generator coverage is that of the
labelled corpus. The formality mechanism rests on a difference in means
rather than a significance-tested effect. The most useful next step follows from
the calibration result: whether a threshold set \emph{without} the held-out
labels --- from the training assignments' score distribution, or from a
small marked sample of the new assignment --- recovers any of the headroom
the oracle threshold exposes. If none of it is recoverable, the case for
detection as an assessment control is weaker than any accuracy figure
suggests. Errors that fall unevenly across students strengthen the case for
assessment designs which do not depend on authorship inference: process
evidence, oral defence, or work situated enough to be difficult to
outsource.

\begin{thebibliography}{99}
\small

\bibitem{qorich2025} Qorich, M. and El Ouazzani, R. (2025) ``Detection of
artificial intelligence-generated essays for academic assessment integrity
using large language models,'' \emph{Expert Systems With Applications},
291, 128405.

\bibitem{crossley2024} Crossley, S. A., Baffour, P., Tian, Y., Franklin,
A., Benner, M. and Boser, U. (2024) ``A large-scale corpus for assessing
written argumentation: PERSUADE 2.0,'' \emph{Assessing Writing}, 61,
100255.

\bibitem{king2023} King, J., Baffour, P., Crossley, S., Holbrook, R. and
Demkin, M. (2023) ``LLM --- Detect AI Generated Text,'' Kaggle competition,
Vanderbilt University and The Learning Agency Lab.

\bibitem{daigt2024} Kleczek, D. (2024) ``DAIGT V2 Train Dataset,'' Kaggle
dataset.
\url{https://www.kaggle.com/datasets/thedrcat/daigt-v2-train-dataset}

\bibitem{ribeiro2016} Ribeiro, M. T., Singh, S. and Guestrin, C. (2016)
```Why should I trust you?'' Explaining the predictions of any classifier,''
in \emph{Proceedings of the 22nd ACM SIGKDD International Conference on
Knowledge Discovery and Data Mining}, pp. 1135--1144.

\bibitem{krishna2023} Krishna, K., Song, Y., Karpinska, M., Wieting, J. and
Iyyer, M. (2023) ``Paraphrasing evades detectors of AI-generated text, but
retrieval is an effective defense,'' in \emph{Proceedings of NeurIPS}, 36.

\bibitem{weberwulff2023} Weber-Wulff, D. et al. (2023) ``Testing of
detection tools for AI-generated text,'' \emph{International Journal for
Educational Integrity}, 19(1), 26.

\bibitem{liang2023} Liang, W., Yuksekgonul, M., Mao, Y., Wu, E. and Zou, J.
(2023) ``GPT detectors are biased against non-native English writers,''
\emph{Patterns}, 4(7), 100779.

\bibitem{sadasivan2023} Sadasivan, V.S., Kumar, A., Balasubramanian, S.,
Wang, W. and Feizi, S. (2023) ``Can AI-generated text be reliably
detected?'' arXiv:2303.11156.

\bibitem{mage2024} Li, Y. et al. (2024) ``MAGE: Machine-generated text
detection in the wild,'' in \emph{Proceedings of ACL}, 36--53.

\bibitem{sentinel2026} Anonymous (2026) ``DeBERTa-Sentinel: toward
transparent and trustworthy detection of AI-generated text.'' Reference
withheld to preserve anonymity; full citation supplied on acceptance.

\bibitem{rag2026} Anonymous (2026) ``Integrating retrieval-augmented
generation with large language model for robust and explainable AI text
detection.'' Reference withheld to preserve anonymity; full citation
supplied on acceptance.

\bibitem{raid2024} Dugan, L., Hwang, A., Trhl{\'i}k, F., Zhu, A., Ludan,
J. M., Xu, H., Ippolito, D. and Callison-Burch, C. (2024) ``RAID: A shared
benchmark for robust evaluation of machine-generated text detectors,'' in
\emph{Proceedings of the 62nd Annual Meeting of the Association for
Computational Linguistics}, Volume 1, pp. 12463--12492.

\end{thebibliography}

\end{document}
